\subsection{A genuinely nonabelian finite obstruction mechanism}

\begin{claim}{B.NA1}{Proved}{A finite norm obstruction with no vanishing normal-step module}
The equivalence of \statementref{B.AB1} cannot be extended to
arbitrary finite groups using only the vanishing of relative
$p$-primary modules at normal quotient steps. There is a
faithful action of a finite nonabelian simple $7'$-group $G$
on a nonzero finite $7$-group $M$ such that the set of vectors
with free norm orbit at every normal quotient is empty,
although the relative module at every nontrivial normal step
is nonzero. This is an exact finite representation example,
not a claimed realization by a Jacobian or a counterexample
to universal B1.
\end{claim}
\begin{proof}
Let $M=\mathbf F_7^5$, and define $G\le\operatorname{GL}(M)$
by the three matrices
\[
 U=\operatorname{diag}(1,4,1,4,4),\qquad
 P=\begin{pmatrix}
 1&0&0&0&0\\
 0&1&0&0&0\\
 0&0&0&0&1\\
 0&0&1&0&0\\
 0&0&0&1&0
 \end{pmatrix},\qquad
 Q=\begin{pmatrix}
 2&4&4&4&4\\
 2&5&4&5&5\\
 2&4&5&5&5\\
 2&5&5&5&4\\
 2&5&5&4&5
 \end{pmatrix}.                                                \tag{B.NA.1}
\]
Here $U^3=P^3=Q^2=1$, and $UP\ne PU$.
The following finite certificate proves both the order and
simplicity needed in the argument, without assuming an
identification with a named finite classical group.

Starting with the identity, close the set of matrices under
right multiplication by $U,P,Q$. Exact arithmetic modulo
seven gives
\[
                        |G|=25920=2^6\cdot3^4\cdot5.
\]
Conjugation by the same generators partitions this complete
matrix list into twenty classes, with sizes
\[
\begin{gathered}
 1,40,40,45,240,270,360,360,480,540,720,720,1440,\\
                   2160,2160,2160,2880,2880,3240,5184.
\end{gathered}                                                   \tag{B.NA.2}
\]
For each nonidentity class representative $a$, impose the
equivalence relation $1\sim a$ on the matrix list and close
it under left and right multiplication by $U,P,Q$. This is
the group congruence whose identity class is the normal
closure of $a$. In all nineteen cases its identity class
has all $25920$ elements. Every nontrivial normal subgroup
contains one of these representatives up to conjugacy, so
$G$ is simple.

The data are fully reproducible. The generating program saves
every matrix in
\path{code/criterion_nonabelian_norm_group.json}. The independent
checker \path{code/criterion_nonabelian_norm_verify.py} reads
the three generators and regenerates the complete group using
a separate general matrix multiplication routine. It checks
the conjugacy partition and the normal closures just described.
Congruence closure uses a finite union-find calculation: after
a new pair is identified, its six images under left and right
generator multiplication are queued. This computes the least
invariant equivalence relation and therefore the asserted
normal closure. The checker also verifies that the common
fixed space of the three generators has dimension zero and
that $\sum_{g\in G}g$ is the zero operator on $M$.

The action is faithful by its definition as a matrix group,
and it is coprime because $7\nmid25920$. Nevertheless
\[
                         |M|=7^5=16807<|G|.
\]
No orbit in $M$ can therefore have size $|G|$. In particular
the norm-freeness condition already fails at the identity
subgroup, where the norm coordinate is the vector itself.

On the other hand, simplicity leaves only the normal
subgroups $1$ and $G$. The old module at the latter level is
$M^G=0$, and the relative module for the sole nontrivial
normal step is the entire nonzero module $M$. Its order is
$7^5$, not a unit. Thus there is no vanishing relative module
at any nontrivial normal step, yet the full finite norm-free
set is empty. The obstruction is the union of nontrivial
stabilizers, rather than one absent abelian character packet.

For an actual cover in B2, the separate task is to determine
whether its Frobenius-fixed Jacobian module has such a
configuration. This example proves that a search for
nonabelian finite obstructions has a mechanism not captured
by the relative class-number-unit sufficient criterion. It
does not supply that geometric realization or settle the
universal assertion.
\end{proof}
